\section{Methodology}
\label{sec:methodology}

Building on our observation that RL training and test-time refinement both leverage execution feedback but have never been combined systematically, we design a factorial experimental framework to measure their interaction and probe the underlying mechanism.

\subsection{Overview}

Our approach consists of three components: (1) a 2$\times$2 factorial design that isolates the Training$\times$Refinement interaction, (2) a mutual information metric $I(F;E)$ that quantifies feedback-edit coupling, and (3) a diversity manipulation infrastructure that enables ablation studies. We describe each component's design rationale and implementation.

\subsection{2$\times$2 Factorial Design}

We cross two factors:

\begin{itemize}
    \item \textbf{Training type:} Cross-entropy (CE) fine-tuning vs.\ RL fine-tuning with execution rewards (REINFORCE)
    \item \textbf{Inference mode:} Single-shot generation vs.\ K-step refinement (Self-Refine protocol, $K=3$)
\end{itemize}

This yields four conditions:

\begin{table}[h]
\centering
\begin{tabular}{lll}
\toprule
Condition & Training & Inference \\
\midrule
CE-Single & Cross-entropy & Single-shot \\
CE-Refine & Cross-entropy & $K=3$ refinement \\
RL-Single & RL + execution & Single-shot \\
RL-Refine & RL + execution & $K=3$ refinement \\
\bottomrule
\end{tabular}
\caption{Four factorial conditions crossing training type with inference mode.}
\label{tab:conditions}
\end{table}

\textbf{Rationale:} Factorial design separates main effects from interaction. The interaction term---whether RL-Refine exceeds what CE-Refine and RL-Single would predict additively---directly tests our hypothesis of superadditivity.

\textbf{Statistical model:} We fit a generalized linear mixed model (GLMM) with logit link:
\begin{equation}
\text{logit}(P(\text{pass})) = \beta_0 + \beta_T \cdot \text{Training} + \beta_R \cdot \text{Refinement} + \beta_{T \times R} \cdot \text{Training} \times \text{Refinement} + u_{\text{problem}}
\end{equation}
where $u_{\text{problem}}$ is a random effect capturing problem-level difficulty. Our success criterion is $\beta_{T \times R} > 0$ with $p < 0.05$ and odds ratio $\geq 1.2$.

\subsection{Mutual Information Metric: $I(F;E)$}

To probe \emph{why} interaction effects occur, we measure structural coupling between feedback and edits.

\textbf{Definition:} $I(F;E)$ is the mutual information between feedback embeddings $F$ and edit embeddings $E$.

\textbf{Implementation:} We use the MINE neural estimator \citep{belghazi2018mine}, which provides the Donsker-Varadhan lower bound on mutual information. We train for 5000 iterations and use 10,000 permutations for significance testing.

\textbf{Control:} We regress out edit length to ensure MI differences are not artifacts of RL producing longer edits.

\subsection{Feedback Diversity Controller}

To test whether diversity is necessary for superadditivity, we manipulate training feedback entropy.

\textbf{Definition:} Feedback diversity is measured via conditional entropy $H(\text{Schema} \mid \text{ErrorClass})$.

\textbf{Achieved separation:} $H_{\text{high}} = 2.0$ bits, $H_{\text{low}} = 0.9$ bits (1.1-bit difference), enabling controlled ablation.
